\documentclass[12pt,a4paper]{article}

% =====================================================================
% PACKAGES
% =====================================================================
\usepackage[utf8]{inputenc}
\usepackage[T1]{fontenc}
\usepackage[german]{babel}
\usepackage{amsmath, amssymb, amsfonts}
\usepackage{geometry}
\usepackage{graphicx}
\usepackage{hyperref}
\usepackage{xcolor}
\usepackage{titlesec}
\usepackage{fancyhdr}
\usepackage{booktabs}
\usepackage{enumitem}
\usepackage{authblk}
\usepackage{setspace}
\usepackage{tcolorbox}
\usepackage{listings}
\usepackage{tikz}
\usetikzlibrary{shapes.geometric, arrows, positioning, fit, backgrounds}

% =====================================================================
% LAYOUT
% =====================================================================
\geometry{left=2.5cm, right=2.5cm, top=3cm, bottom=3cm}
\onehalfspacing

% =====================================================================
% COLORS & LINKS
% =====================================================================
\definecolor{darkblue}{RGB}{0,51,102}
\definecolor{codegreen}{RGB}{0,128,0}
\definecolor{codegray}{RGB}{128,128,128}
\definecolor{codepurple}{RGB}{128,0,128}
\definecolor{backcolour}{RGB}{245,245,245}
\hypersetup{
    colorlinks=true,
    linkcolor=darkblue,
    citecolor=darkblue,
    urlcolor=darkblue
}

% =====================================================================
% CODE LISTINGS
% =====================================================================
\lstdefinestyle{pythonstyle}{
    backgroundcolor=\color{backcolour},
    commentstyle=\color{codegreen},
    keywordstyle=\color{codepurple},
    numberstyle=\tiny\color{codegray},
    stringstyle=\color{codepurple},
    basicstyle=\ttfamily\footnotesize,
    breakatwhitespace=false,
    breaklines=true,
    captionpos=b,
    keepspaces=true,
    numbers=left,
    numbersep=5pt,
    showspaces=false,
    showstringspaces=false,
    showtabs=false,
    tabsize=2,
    language=Python
}
\lstset{style=pythonstyle}

% =====================================================================
% SECTION FORMATTING
% =====================================================================
\titleformat{\section}{\Large\bfseries\color{darkblue}}{\thesection}{1em}{}
\titleformat{\subsection}{\large\bfseries\color{darkblue!80}}{\thesubsection}{1em}{}

% =====================================================================
% HEADER & FOOTER
% =====================================================================
\pagestyle{fancy}
\fancyhf{}
\fancyhead[R]{\small\color{gray}BEATRIX Drei-Phasen-Architektur}
\fancyfoot[C]{\thepage}
\renewcommand{\headrulewidth}{0.4pt}

% =====================================================================
% DOCUMENT INFO
% =====================================================================
\title{
{\LARGE\bfseries Technisches Dokument}\\[0.4cm]
{\Large Die Drei-Phasen-Architektur von BEATRIX}\\[0.2cm]
{\large LLM-Orchestrierung für verhaltensökonomische Modellierung}
}

\author[1]{Gerhard Fehr}
\author[1]{Andrea Fehr}
\affil[1]{FehrAdvice \& Partners AG, Zürich}
\date{Dezember 2025 | Version 1.0}

% =====================================================================
% DOCUMENT
% =====================================================================
\begin{document}

\maketitle

\begin{abstract}
Dieses technische Dokument beschreibt die Drei-Phasen-Architektur von BEATRIX, die eine nahtlose Integration von Large Language Models (LLMs) mit dem Behavioral Change Model (BCM 2.2) ermöglicht.
Die Architektur löst ein fundamentales Problem der Pre-LLM-Ära: Das BCM konnte berechnet werden, war aber für Nicht-Experten schwer zugänglich.
BEATRIX orchestriert LLMs und deterministischen Code in drei Phasen: (1) LLM-basierte Extraktion von BCM-Parametern aus natürlicher Sprache, (2) deterministische BCM-Berechnung via Code Interpreter, (3) LLM-basierte Interpretation der Ergebnisse.
Wir dokumentieren die historische Entwicklung der enabling technologies (möglich seit Juni 2023, production-ready seit August 2024), die wissenschaftlichen Grundlagen aus zehn Forschungsbereichen, und drei technische Implementierungsvarianten mit konkreten Code-Beispielen.
\end{abstract}

\tableofcontents
\newpage

% =====================================================================
\section{Einleitung: Das Problem und die Lösung}
% =====================================================================

\subsection{Das Pre-LLM Problem}

Das Behavioral Change Model (BCM) existiert seit 2023 als mathematisches Framework zur Modellierung menschlichen Verhaltens.
Die Kernformel für den effektiven Nutzen lautet:

\begin{equation}
U_{\text{eff}} = (1-\lambda) \cdot U_{\text{INU}} + \lambda \cdot U_{\text{KNU}} + U_{\text{IDN}}(a, \sigma, \phi, \tilde{\mu}; F, s, m)
\end{equation}

Das BCM konnte immer \textbf{berechnet} werden -- gegeben die richtigen Parameter, liefert die Mathematik deterministische Ergebnisse.
Das Problem lag an zwei Stellen:

\begin{enumerate}
    \item \textbf{Input-Problem}: Wie kommt man von einer Persona-Beschreibung (``Anna ist 45, Lehrerin, umweltbewusst...'') zu den Modellparametern ($\lambda = 2.2$, $\alpha = 0.88$, ...)?

    \item \textbf{Output-Problem}: Wie erklärt man einem Nicht-Experten, was ``$U_{\text{eff}} = -0.34$ mit Sensitivität $\partial U / \partial \lambda = 0.6$'' bedeutet?
\end{enumerate}

Vor LLMs erforderten beide Schritte \textbf{menschliche Experten} -- Verhaltensökonomen, die Parameter schätzen und Ergebnisse interpretieren konnten.
Dies machte BCM-Analysen langsam, teuer und nicht skalierbar.

\subsection{Die BEATRIX-Lösung: Drei Phasen}

BEATRIX löst dieses Problem durch eine orchestrierte Drei-Phasen-Architektur:

\begin{figure}[h]
\centering
\begin{tikzpicture}[
    node distance=1.5cm,
    phase/.style={rectangle, draw, rounded corners, minimum width=3.5cm, minimum height=2cm, align=center},
    arrow/.style={->, thick, >=stealth}
]

% Phases
\node[phase, fill=blue!15] (p1) at (0, 0) {\textbf{Phase 1}\\LLM\\[0.2em]\footnotesize Input-Parsing};
\node[phase, fill=green!15] (p2) at (5, 0) {\textbf{Phase 2}\\Code Interpreter\\[0.2em]\footnotesize BCM-Berechnung};
\node[phase, fill=orange!15] (p3) at (10, 0) {\textbf{Phase 3}\\LLM\\[0.2em]\footnotesize Interpretation};

% Input/Output
\node[above=0.8cm of p1, align=center] (input) {\footnotesize ``Anna ist 45,\\Lehrerin, ...''};
\node[above=0.8cm of p2, align=center] (params) {\footnotesize $\lambda=2.2$, $\alpha=0.88$\\$w_F=0.35$, ...};
\node[above=0.8cm of p3, align=center] (results) {\footnotesize Churn: 4.2\%\\CI: [2.8\%, 5.9\%]};
\node[below=0.8cm of p3, align=center] (output) {\footnotesize ``Bei 5\% Preiserhöhung\\erwarten wir...''};

% Arrows
\draw[arrow] (input) -- (p1);
\draw[arrow] (p1) -- node[above, font=\footnotesize] {JSON} (p2);
\draw[arrow] (p2) -- node[above, font=\footnotesize] {Dict} (p3);
\draw[arrow] (p3) -- (output);

% Labels below
\node[below=0.5cm of p1, font=\footnotesize\itshape, text=gray] {Probabilistisch};
\node[below=0.5cm of p2, font=\footnotesize\itshape, text=gray] {Deterministisch};
\node[below=0.5cm of p3, font=\footnotesize\itshape, text=gray] {Probabilistisch};

\end{tikzpicture}
\caption{Die Drei-Phasen-Architektur von BEATRIX. LLMs übernehmen die ``Übersetzung'' zwischen natürlicher Sprache und Modellparametern; die eigentliche Berechnung erfolgt deterministisch.}
\label{fig:three-phases}
\end{figure}

\begin{tcolorbox}[colback=blue!5!white, colframe=darkblue, title=Kernprinzip]
Das LLM ist \textbf{kein Taschenrechner} -- es führt keine BCM-Berechnungen durch.\\
Das LLM ist ein \textbf{Übersetzer}:
\begin{itemize}
    \item Phase 1: Mensch $\rightarrow$ Modell (natürliche Sprache $\rightarrow$ Parameter)
    \item Phase 3: Modell $\rightarrow$ Mensch (Ergebnisse $\rightarrow$ Narrative)
\end{itemize}
Die \textbf{Mathematik bleibt deterministisch} im Code Interpreter (Phase 2).
\end{tcolorbox}

% =====================================================================
\section{Historische Entwicklung: Seit wann ist das möglich?}
% =====================================================================

Die Drei-Phasen-Architektur wurde erst durch eine Konvergenz mehrerer technologischer Entwicklungen möglich.
Dieser Abschnitt dokumentiert die relevanten Meilensteine.

\subsection{Timeline der Enabling Technologies}

\begin{table}[h]
\centering
\caption{Meilensteine der LLM-Entwicklung für BEATRIX}
\small
\begin{tabular}{@{}llp{7cm}@{}}
\toprule
\textbf{Datum} & \textbf{Entwicklung} & \textbf{Bedeutung für BEATRIX} \\
\midrule
Jun 2017 & Transformer (Vaswani et al.) & Grundarchitektur für alle modernen LLMs \\
Jun 2018 & GPT-1 (117M Parameter) & Erstes generatives Pre-Training \\
Jun 2020 & GPT-3 (175B Parameter) & In-Context Learning entdeckt \\
\midrule
\textbf{Mär 2022} & \textbf{InstructGPT / ChatGPT} & \textbf{Instruction Following wird zuverlässig} \\
& & $\rightarrow$ Phase 1 (Input) wird möglich \\
Dez 2022 & Chain-of-Thought (Wei et al.) & Reasoning-Fähigkeiten verbessert \\
\midrule
Mär 2023 & GPT-4 & Deutlich verbesserte Reasoning-Qualität \\
\textbf{Jun 2023} & \textbf{Function Calling (OpenAI)} & \textbf{Strukturierte Outputs werden zuverlässig} \\
& & $\rightarrow$ JSON-Parameter-Extraktion möglich \\
\textbf{Jul 2023} & \textbf{Code Interpreter (ChatGPT)} & \textbf{Python-Ausführung in Sandbox} \\
& & $\rightarrow$ Phase 2 wird nahtlos integrierbar \\
Jul 2023 & Claude 2 + Tool Use & Alternative zu OpenAI, 100k Kontext \\
\midrule
Mär 2024 & Claude 3 (Opus, Sonnet, Haiku) & 200k Kontext, robustere Tool Use \\
\textbf{Aug 2024} & \textbf{Structured Outputs (OpenAI)} & \textbf{JSON Schema Enforcement} \\
& & $\rightarrow$ \textbf{Garantierte} Struktur \\
Okt 2024 & Claude 3.5 Sonnet & Verbesserte Reasoning-Fähigkeiten \\
\bottomrule
\end{tabular}
\end{table}

\subsection{Kritische Schwellenwerte}

\begin{tcolorbox}[colback=yellow!5!white, colframe=yellow!75!black, title=Zeitliche Einordnung]
\textbf{Frühester Zeitpunkt} (Proof of Concept): \textbf{Juni 2023}\\
Mit Function Calling und Code Interpreter waren alle drei Phasen grundsätzlich umsetzbar.

\textbf{Production-Ready Zeitpunkt}: \textbf{August 2024}\\
Mit Structured Outputs (JSON Schema Enforcement) werden Parameter-Extraktionen \textit{garantiert} strukturiert -- nicht nur ``wahrscheinlich''.

\textbf{BEATRIX-Architektur ist erst seit $\approx$18 Monaten praktisch umsetzbar.}
\end{tcolorbox}

\subsection{Warum nicht früher?}

Vor den genannten Entwicklungen fehlten kritische Fähigkeiten:

\begin{table}[h]
\centering
\caption{Fehlende Fähigkeiten vor 2023}
\begin{tabular}{@{}lll@{}}
\toprule
\textbf{Phase} & \textbf{Benötigte Fähigkeit} & \textbf{Verfügbar seit} \\
\midrule
Phase 1 & Zuverlässiges Instruction Following & März 2022 \\
Phase 1 & Strukturierte JSON-Outputs & Juni 2023 \\
Phase 1 & Garantierte Schema-Konformität & August 2024 \\
Phase 2 & Code-Ausführung in LLM-Kontext & Juli 2023 \\
Phase 3 & Lange Kontexte (>50k Tokens) & Juli 2023 \\
Phase 3 & Grounded Generation (weniger Halluzination) & 2024 \\
\bottomrule
\end{tabular}
\end{table}

% =====================================================================
\section{Wissenschaftliche Grundlagen}
% =====================================================================

Die Drei-Phasen-Architektur stützt sich auf Erkenntnisse aus zehn Forschungsbereichen.

\subsection{Phase 1: Natural Language $\rightarrow$ Parameter}

\subsubsection{Information Extraction (IE)}

Die Extraktion strukturierter Information aus unstrukturiertem Text ist ein klassisches NLP-Problem:

\begin{itemize}
    \item \textbf{Named Entity Recognition (NER)}: Identifikation von Entitäten (Personen, Orte, Organisationen)
    \item \textbf{Relation Extraction}: Identifikation von Beziehungen zwischen Entitäten
    \item \textbf{Slot Filling}: Extraktion von Attributwerten für vordefinierte Slots
\end{itemize}

BEATRIX nutzt LLMs für eine erweiterte Form von Slot Filling: Die ``Slots'' sind BCM-Parameter, und die ``Werte'' sind numerische Schätzungen mit Reasoning.

\textbf{Schlüsselliteratur}: Jurafsky \& Martin (2024), Kapitel 17.

\subsubsection{Semantic Parsing}

Semantic Parsing transformiert natürliche Sprache in formale Repräsentationen:

\begin{quote}
``Anna ist preisbewusst'' $\rightarrow$ \texttt{w\_financial: 0.4, lambda: 2.3}
\end{quote}

Klassische Ansätze nutzten Grammatiken und regelbasierte Systeme.
LLMs ermöglichen \textbf{zero-shot} Semantic Parsing durch In-Context Learning.

\textbf{Schlüsselliteratur}: Liang (2016), ``Learning Executable Semantic Parsers''.

\subsubsection{In-Context Learning}

Brown et al. (2020) zeigten, dass LLMs neue Aufgaben ohne Finetuning erlernen können, allein durch Beispiele im Prompt:

\begin{lstlisting}
# Few-Shot Prompting für Parameter-Extraktion
Beispiel 1:
Input: "Max ist 25, Student, technikaffin"
Output: {"lambda": 1.8, "w_digital": 0.5, ...}

Beispiel 2:
Input: "Petra ist 60, Rentnerin, konservativ"
Output: {"lambda": 2.4, "w_financial": 0.3, ...}

Jetzt du:
Input: "Anna ist 45, Lehrerin, umweltbewusst"
Output: ?
\end{lstlisting}

\textbf{Schlüsselliteratur}: Brown et al. (2020), GPT-3 Paper; Min et al. (2022), ``Rethinking In-Context Learning''.

\subsubsection{Instruction Tuning}

RLHF (Reinforcement Learning from Human Feedback) transformiert Base-Modelle in Instruction-Following-Assistenten:

\begin{itemize}
    \item \textbf{Supervised Fine-Tuning (SFT)}: Training auf Instruction-Response-Paaren
    \item \textbf{Reward Modeling}: Lernen menschlicher Präferenzen
    \item \textbf{PPO/DPO}: Optimierung auf Reward-Signal
\end{itemize}

Ohne Instruction Tuning würden LLMs Text fortsetzen, aber nicht strukturiert antworten.

\textbf{Schlüsselliteratur}: Ouyang et al. (2022), InstructGPT; Bai et al. (2022), Constitutional AI.

\subsection{Phase 2: Parameter $\rightarrow$ Berechnung}

\subsubsection{Behavioral Economics Models}

Das BCM integriert Erkenntnisse aus vier Jahrzehnten verhaltensökonomischer Forschung:

\begin{itemize}
    \item \textbf{Prospect Theory}: Referenzpunkte, Loss Aversion, Probability Weighting (Kahneman \& Tversky, 1979)
    \item \textbf{Reference-Dependent Utility}: Expectations as Reference Points (K\H{o}szegi \& Rabin, 2006)
    \item \textbf{Social Preferences}: Inequity Aversion, Reciprocity (Fehr \& Schmidt, 1999)
    \item \textbf{Identity Utility}: Social Identity, Group Membership (Akerlof \& Kranton, 2000)
\end{itemize}

\textbf{Schlüsselliteratur}: BCM 2.2 Technical Report (Fehr \& Fehr, 2024).

\subsubsection{Uncertainty Quantification}

BCM-Berechnungen müssen Unsicherheit propagieren:

\begin{itemize}
    \item \textbf{Monte Carlo Simulation}: Sampling aus Parameterverteilungen
    \item \textbf{Sensitivity Analysis}: $\partial U / \partial \theta$ für jeden Parameter
    \item \textbf{Confidence Intervals}: Quantile der Output-Verteilung
\end{itemize}

\textbf{Schlüsselliteratur}: Saltelli et al. (2008), ``Global Sensitivity Analysis''.

\subsubsection{Tool-Augmented LLMs}

LLMs können externe Tools aufrufen, um ihre Limitierungen zu überwinden:

\begin{itemize}
    \item \textbf{Toolformer} (Schick et al., 2023): LLMs lernen, wann Tools nützlich sind
    \item \textbf{PAL} (Gao et al., 2023): Program-Aided Language Models für mathematische Aufgaben
    \item \textbf{Code Interpreter}: Native Integration von Python-Ausführung
\end{itemize}

\textbf{Kerninsight}: LLMs sind schlecht in exakter Arithmetik, aber gut darin zu erkennen, \textit{wann} Arithmetik nötig ist.

\textbf{Schlüsselliteratur}: Schick et al. (2023), ``Toolformer''; Gao et al. (2023), ``PAL''.

\subsection{Phase 3: Ergebnis $\rightarrow$ Narrative}

\subsubsection{Natural Language Generation (NLG)}

Data-to-Text Generation transformiert strukturierte Daten in natürliche Sprache:

\begin{quote}
\texttt{\{churn: 0.042, ci: [0.028, 0.059]\}} $\rightarrow$ ``Wir erwarten etwa 4\% Churn, mit einer Spanne von 3-6\%.''
\end{quote}

Klassische NLG-Systeme waren regelbasiert und domänenspezifisch.
LLMs ermöglichen flexible, kontextsensitive Generierung.

\textbf{Schlüsselliteratur}: Reiter \& Dale (2000), ``Building NLG Systems''; Gatt \& Krahmer (2018), ``Survey of NLG''.

\subsubsection{Explainable AI (XAI)}

Erklärungen müssen für den Empfänger verständlich sein:

\begin{itemize}
    \item \textbf{Contrastive Explanations}: ``Warum 4\% und nicht 2\%?''
    \item \textbf{Counterfactual Explanations}: ``Was müsste anders sein für 2\%?''
    \item \textbf{Feature Attribution}: ``$\lambda$ erklärt 60\% der Unsicherheit''
\end{itemize}

\textbf{Schlüsselliteratur}: Miller (2019), ``Explanation in AI''; Molnar (2022), ``Interpretable ML''.

\subsubsection{Grounded Generation}

LLM-Outputs müssen auf den tatsächlichen BCM-Ergebnissen basieren, nicht halluziniert sein:

\begin{itemize}
    \item \textbf{Factual Consistency}: Aussagen müssen mit Input konsistent sein
    \item \textbf{Source Attribution}: ``Laut BCM-Berechnung...''
    \item \textbf{Uncertainty Communication}: Konfidenzintervalle müssen kommuniziert werden
\end{itemize}

\textbf{Schlüsselliteratur}: Ji et al. (2023), ``Survey of Hallucination in NLG''.

% =====================================================================
\section{Technische Implementierung}
% =====================================================================

Wir beschreiben drei Implementierungsvarianten mit steigender Komplexität und Kontrolle.

\subsection{Variante A: Monolithisch (Prototyping)}

Die einfachste Variante nutzt ChatGPT oder Claude mit Code Interpreter direkt:

\begin{figure}[h]
\centering
\begin{tikzpicture}[
    node distance=0.8cm,
    box/.style={rectangle, draw, rounded corners, minimum width=8cm, minimum height=1cm, align=center}
]

\node[box, fill=gray!10] (llm) at (0, 0) {LLM (ChatGPT/Claude mit Code Interpreter)};
\node[above=0.5cm of llm] (input) {User Input};
\node[below=0.5cm of llm] (output) {Interpretation};

\draw[->] (input) -- (llm);
\draw[->] (llm) -- (output);

\node[right=0.5cm of llm, align=left, font=\footnotesize] {
    1. LLM parst Input\\
    2. LLM schreibt BCM-Code\\
    3. Code wird ausgeführt\\
    4. LLM interpretiert Output
};

\end{tikzpicture}
\caption{Monolithische Architektur: Alles in einer LLM-Session.}
\end{figure}

\textbf{Vorteile}:
\begin{itemize}
    \item Sehr schnell zu implementieren
    \item Gut für Exploration und Demos
    \item Keine Backend-Infrastruktur nötig
\end{itemize}

\textbf{Nachteile}:
\begin{itemize}
    \item BCM-Code wird jedes Mal neu generiert (Fehleranfällig)
    \item Keine Validierung der Parameter
    \item Nicht reproduzierbar
    \item Schwer zu debuggen
\end{itemize}

\textbf{Eignung}: Proof of Concept, Demos, Exploration.

\subsection{Variante B: API-Orchestriert (Production)}

Die Production-Variante trennt die drei Phasen explizit:

\begin{figure}[h]
\centering
\begin{tikzpicture}[
    node distance=1cm,
    phase/.style={rectangle, draw, rounded corners, minimum width=3cm, minimum height=1.5cm, align=center, font=\small},
    arrow/.style={->, thick}
]

\node[phase, fill=blue!10] (p1) at (0, 0) {Phase 1\\LLM API};
\node[phase, fill=green!10] (p2) at (4.5, 0) {Phase 2\\BCM Engine};
\node[phase, fill=orange!10] (p3) at (9, 0) {Phase 3\\LLM API};

\node[above=0.8cm of p2, draw, dashed, minimum width=10cm, minimum height=0.8cm] (orch) {BEATRIX Orchestrator (Python)};

\draw[arrow] (p1) -- node[above, font=\footnotesize] {JSON} (p2);
\draw[arrow] (p2) -- node[above, font=\footnotesize] {Dict} (p3);

\draw[arrow, dashed] (orch) -- (p1);
\draw[arrow, dashed] (orch) -- (p2);
\draw[arrow, dashed] (orch) -- (p3);

\end{tikzpicture}
\caption{API-orchestrierte Architektur: Explizite Trennung der Phasen.}
\end{figure}

\subsubsection{Phase 1: Parameter-Extraktion}

\begin{lstlisting}[caption={Phase 1: LLM-basierte Parameter-Extraktion}]
from anthropic import Anthropic
from pydantic import BaseModel, Field

class BCMParameters(BaseModel):
    """Pydantic Schema fuer strukturierte Outputs"""
    lambda_loss_aversion: float = Field(
        ge=1.0, le=3.0,
        description="Loss Aversion Koeffizient"
    )
    alpha: float = Field(
        ge=0.0, le=1.0,
        description="Value Function Parameter (Gewinne)"
    )
    beta: float = Field(
        ge=0.0, le=1.0,
        description="Value Function Parameter (Verluste)"
    )
    w_financial: float = Field(ge=0.0, le=1.0)
    w_emotional: float = Field(ge=0.0, le=1.0)
    w_ecological: float = Field(ge=0.0, le=1.0)
    reasoning: str = Field(
        description="Chain-of-Thought Begruendung"
    )

def extract_parameters(persona: str) -> BCMParameters:
    client = Anthropic()

    response = client.messages.create(
        model="claude-sonnet-4-20250514",
        max_tokens=4096,
        system="""Du bist ein BCM-Parameter-Extraktor.

        Base Rates (Anker):
        - lambda (Loss Aversion): 2.0 [1.0-3.0]
        - alpha (Gains): 0.88 [0.7-1.0]
        - beta (Losses): 0.88 [0.7-1.0]

        Dokumentiere dein Reasoning.""",
        messages=[{
            "role": "user",
            "content": f"Extrahiere BCM-Parameter:\n{persona}"
        }],
        tools=[{
            "name": "extract_bcm",
            "input_schema": BCMParameters.model_json_schema()
        }],
        tool_choice={"type": "tool", "name": "extract_bcm"}
    )

    return BCMParameters(**response.content[0].input)
\end{lstlisting}

\subsubsection{Phase 2: BCM-Berechnung}

\begin{lstlisting}[caption={Phase 2: Deterministische BCM-Berechnung}]
import numpy as np
from typing import Dict

def calculate_bcm(
    params: BCMParameters,
    scenario: Dict
) -> Dict:
    """Deterministische BCM 2.2 Berechnung"""

    # Prospect Theory Value Function
    def value_function(x, alpha, beta, lambda_):
        if x >= 0:
            return x ** alpha
        else:
            return -lambda_ * ((-x) ** beta)

    # Berechne Utilities
    delta_price = scenario.get("price_change", 0.05)

    # Individual Utility (vereinfacht)
    U_INU = value_function(
        -delta_price,  # Preiserhoehung = Verlust
        params.alpha,
        params.beta,
        params.lambda_loss_aversion
    )

    # FEPSDE-gewichteter Utility
    fepsde_weights = np.array([
        params.w_financial,
        params.w_emotional,
        0.1,  # Physical (default)
        0.15, # Social (default)
        0.1,  # Digital (default)
        params.w_ecological
    ])
    fepsde_weights = fepsde_weights / fepsde_weights.sum()

    # Monte Carlo fuer Unsicherheit
    n_samples = 10000
    samples = []

    for _ in range(n_samples):
        # Sample Parameter mit Unsicherheit
        l = np.random.normal(
            params.lambda_loss_aversion, 0.2
        )
        l = np.clip(l, 1.0, 3.0)

        u = value_function(-delta_price, params.alpha,
                          params.beta, l)
        samples.append(u)

    samples = np.array(samples)

    # Churn-Wahrscheinlichkeit (Logit-Transform)
    churn_prob = 1 / (1 + np.exp(-samples * 10))

    return {
        "U_eff": float(np.mean(samples)),
        "churn_probability": float(np.mean(churn_prob)),
        "confidence_interval": [
            float(np.percentile(churn_prob, 5)),
            float(np.percentile(churn_prob, 95))
        ],
        "sensitivity": {
            "lambda": float(np.std(churn_prob) / 0.2),
        }
    }
\end{lstlisting}

\subsubsection{Phase 3: Interpretation}

\begin{lstlisting}[caption={Phase 3: LLM-basierte Interpretation}]
def interpret_results(
    params: BCMParameters,
    results: Dict,
    persona: str,
    scenario: Dict
) -> str:
    """LLM interpretiert BCM-Ergebnisse"""

    client = Anthropic()

    response = client.messages.create(
        model="claude-sonnet-4-20250514",
        max_tokens=2048,
        system="""Du bist ein Verhaltensoekonimie-Berater.

        Regeln:
        - Erklaere Ergebnisse verstaendlich
        - Vermeide Fachjargon
        - Gib konkrete Handlungsempfehlungen
        - Kommuniziere Unsicherheit ehrlich
        - Beziehe dich auf die Persona""",
        messages=[{
            "role": "user",
            "content": f"""
Persona: {persona}

Szenario: {scenario}

Extrahierte Parameter:
- Loss Aversion (lambda): {params.lambda_loss_aversion}
- Reasoning: {params.reasoning}

BCM-Ergebnisse:
- Erwarteter Churn: {results['churn_probability']:.1%}
- Konfidenzintervall: [{results['confidence_interval'][0]:.1%},
                       {results['confidence_interval'][1]:.1%}]
- Sensitivitaet Lambda: {results['sensitivity']['lambda']:.2f}

Bitte interpretiere diese Ergebnisse.
            """
        }]
    )

    return response.content[0].text
\end{lstlisting}

\subsubsection{Orchestrator}

\begin{lstlisting}[caption={BEATRIX Orchestrator: Integration der drei Phasen}]
def beatrix_analyze(
    persona: str,
    scenario: Dict
) -> Dict:
    """Hauptfunktion: Orchestriert alle drei Phasen"""

    # Phase 1: Extract
    print("Phase 1: Extrahiere Parameter...")
    params = extract_parameters(persona)

    # Validierung
    assert 1.0 <= params.lambda_loss_aversion <= 3.0
    assert params.w_financial + params.w_emotional + \
           params.w_ecological <= 1.0

    # Phase 2: Calculate
    print("Phase 2: Berechne BCM...")
    results = calculate_bcm(params, scenario)

    # Phase 3: Interpret
    print("Phase 3: Interpretiere Ergebnisse...")
    interpretation = interpret_results(
        params, results, persona, scenario
    )

    return {
        "parameters": params.model_dump(),
        "results": results,
        "interpretation": interpretation,
        "audit_trail": {
            "phase1_model": "claude-sonnet-4-20250514",
            "phase1_reasoning": params.reasoning,
            "phase2_deterministic": True,
            "phase3_model": "claude-sonnet-4-20250514",
            "timestamp": datetime.now().isoformat()
        }
    }

# Beispielaufruf
if __name__ == "__main__":
    result = beatrix_analyze(
        persona="Anna ist 45, Lehrerin, umweltbewusst, "
                "preissensitiv, hat zwei Kinder",
        scenario={"price_change": 0.05}
    )
    print(result["interpretation"])
\end{lstlisting}

\subsection{Variante C: Agentic (Autonomer Multi-Step)}

Die fortgeschrittenste Variante nutzt einen autonomen Agenten:

\begin{figure}[h]
\centering
\begin{tikzpicture}[
    node distance=1cm,
    brain/.style={ellipse, draw, minimum width=3cm, minimum height=1.5cm, align=center, fill=purple!10},
    tool/.style={rectangle, draw, rounded corners, minimum width=2.2cm, minimum height=1cm, align=center, fill=gray!10, font=\small}
]

\node[brain] (llm) at (0, 0) {LLM\\(Agent Brain)};

\node[tool] (t1) at (-4, -2.5) {extract\_\\parameters};
\node[tool] (t2) at (0, -2.5) {calculate\_\\bcm};
\node[tool] (t3) at (4, -2.5) {explain\_\\results};

\draw[->] (llm) -- (t1);
\draw[->] (llm) -- (t2);
\draw[->] (llm) -- (t3);

\draw[->, dashed] (t1) to[bend left] (llm);
\draw[->, dashed] (t2) -- (llm);
\draw[->, dashed] (t3) to[bend right] (llm);

\node[below=0.3cm of t2, font=\footnotesize\itshape] {Tools geben Ergebnisse zurück};
\node[above=0.5cm of llm, font=\footnotesize\itshape] {Agent entscheidet nächsten Schritt};

\end{tikzpicture}
\caption{Agentic Architektur: LLM entscheidet autonom über Tool-Aufrufe.}
\end{figure}

\textbf{Vorteile}:
\begin{itemize}
    \item Kann Rückfragen stellen (``Welche Branche?'')
    \item Kann iterieren (``Parameter unplausibel, nochmal prüfen'')
    \item Flexibel bei unerwarteten Inputs
\end{itemize}

\textbf{Nachteile}:
\begin{itemize}
    \item Weniger vorhersehbar
    \item Mehr LLM-Calls (höhere Kosten)
    \item Schwerer zu testen
\end{itemize}

\textbf{Eignung}: Komplexe Szenarien, interaktive Beratung, explorative Analysen.

% =====================================================================
\section{Validierung und Qualitätssicherung}
% =====================================================================

\subsection{Phase 1: Parameter-Validierung}

\begin{table}[h]
\centering
\caption{Validierungsregeln für BCM-Parameter}
\begin{tabular}{@{}llll@{}}
\toprule
\textbf{Parameter} & \textbf{Range} & \textbf{Base Rate} & \textbf{Plausibilitätscheck} \\
\midrule
$\lambda$ & [1.0, 3.0] & 2.0 & $|\lambda - 2.0| < 1.0$ typisch \\
$\alpha$ & [0.0, 1.0] & 0.88 & $\alpha > 0.7$ typisch \\
$\beta$ & [0.0, 1.0] & 0.88 & $\beta \geq \alpha$ (Loss Aversion) \\
$\sum w_i$ & = 1.0 & -- & Normalisierung erforderlich \\
\bottomrule
\end{tabular}
\end{table}

\subsection{Phase 2: Berechnungs-Validierung}

\begin{itemize}
    \item \textbf{Unit Tests}: Bekannte Input-Output-Paare
    \item \textbf{Constraint Checks}: $\beta_{FS} \leq \alpha$, etc.
    \item \textbf{Boundary Tests}: Extremwerte der Parameter
    \item \textbf{Determinismus-Test}: Gleicher Input $\rightarrow$ gleiches Output
\end{itemize}

\subsection{Phase 3: Output-Validierung}

\begin{itemize}
    \item \textbf{Factual Consistency}: Zahlen im Text = Zahlen im Result
    \item \textbf{Uncertainty Communication}: CI muss erwähnt werden
    \item \textbf{No Hallucination}: Keine erfundenen Fakten
\end{itemize}

% =====================================================================
\section{Zusammenfassung und Ausblick}
% =====================================================================

\subsection{Kernbeiträge}

Dieses Dokument hat die Drei-Phasen-Architektur von BEATRIX dokumentiert:

\begin{enumerate}
    \item \textbf{Historische Einordnung}: Die Architektur ist erst seit Juni 2023 grundsätzlich möglich, seit August 2024 production-ready.

    \item \textbf{Wissenschaftliche Fundierung}: Zehn Forschungsbereiche liefern die theoretischen Grundlagen.

    \item \textbf{Technische Implementierung}: Drei Varianten mit steigender Komplexität und Kontrolle.

    \item \textbf{Kernprinzip}: LLMs als Übersetzer, nicht als Rechner. Die BCM-Mathematik bleibt deterministisch.
\end{enumerate}

\subsection{Ausblick}

Zukünftige Entwicklungen:

\begin{itemize}
    \item \textbf{Finetuning}: Spezialisierte LLMs für BCM-Parameter-Extraktion
    \item \textbf{Multi-Modal}: Integration von Bildern, Audio (Interview-Transkripte)
    \item \textbf{Real-Time}: Streaming-Outputs für interaktive Anwendungen
    \item \textbf{Feedback Loop}: Automatische Kalibrierung basierend auf Outcomes
\end{itemize}

\begin{tcolorbox}[colback=green!5!white, colframe=green!75!black, title=Schlussfolgerung]
Die Drei-Phasen-Architektur von BEATRIX transformiert das BCM von einem Experten-Tool zu einem zugänglichen System:

\textbf{Pre-BEATRIX}: Experte $\rightarrow$ Parameter $\rightarrow$ Excel $\rightarrow$ Experte $\rightarrow$ Bericht

\textbf{Mit BEATRIX}: User $\rightarrow$ LLM $\rightarrow$ BCM $\rightarrow$ LLM $\rightarrow$ Verständlicher Output

Der Schlüssel ist die \textbf{kluge Arbeitsteilung}: LLMs für Sprache, Code für Mathematik.
\end{tcolorbox}

% =====================================================================
\section*{Literaturverzeichnis}
% =====================================================================

\begin{thebibliography}{99}

\bibitem[Akerlof \& Kranton, 2000]{akerlof2000}
Akerlof, G. A., \& Kranton, R. E. (2000). Economics and Identity. \textit{Quarterly Journal of Economics, 115}(3), 715--753.

\bibitem[Bai et al., 2022]{bai2022}
Bai, Y., et al. (2022). Constitutional AI: Harmlessness from AI Feedback. arXiv:2212.08073.

\bibitem[Brown et al., 2020]{brown2020}
Brown, T., et al. (2020). Language Models are Few-Shot Learners. \textit{NeurIPS 2020}.

\bibitem[Fehr \& Fehr, 2024]{fehr2024}
Fehr, G., \& Fehr, A. (2024). Behavioral Change Modelling (BCM 2.2). FehrAdvice Technical Report.

\bibitem[Fehr \& Schmidt, 1999]{fehr1999}
Fehr, E., \& Schmidt, K. M. (1999). A Theory of Fairness, Competition, and Cooperation. \textit{Quarterly Journal of Economics, 114}(3), 817--868.

\bibitem[Gao et al., 2023]{gao2023}
Gao, L., et al. (2023). PAL: Program-Aided Language Models. \textit{ICML 2023}.

\bibitem[Gatt \& Krahmer, 2018]{gatt2018}
Gatt, A., \& Krahmer, E. (2018). Survey of the State of the Art in Natural Language Generation. \textit{JAIR, 61}, 65--170.

\bibitem[Ji et al., 2023]{ji2023}
Ji, Z., et al. (2023). Survey of Hallucination in Natural Language Generation. \textit{ACM Computing Surveys, 55}(12).

\bibitem[Jurafsky \& Martin, 2024]{jurafsky2024}
Jurafsky, D., \& Martin, J. H. (2024). \textit{Speech and Language Processing} (3rd ed.). Online Draft.

\bibitem[Kahneman \& Tversky, 1979]{kahneman1979}
Kahneman, D., \& Tversky, A. (1979). Prospect Theory: An Analysis of Decision under Risk. \textit{Econometrica, 47}(2), 263--291.

\bibitem[K\H{o}szegi \& Rabin, 2006]{koszegi2006}
K\H{o}szegi, B., \& Rabin, M. (2006). A Model of Reference-Dependent Preferences. \textit{Quarterly Journal of Economics, 121}(4), 1133--1165.

\bibitem[Liang, 2016]{liang2016}
Liang, P. (2016). Learning Executable Semantic Parsers for Natural Language Understanding. \textit{CACM, 59}(9), 68--76.

\bibitem[Miller, 2019]{miller2019}
Miller, T. (2019). Explanation in Artificial Intelligence: Insights from the Social Sciences. \textit{Artificial Intelligence, 267}, 1--38.

\bibitem[Min et al., 2022]{min2022}
Min, S., et al. (2022). Rethinking the Role of Demonstrations: What Makes In-Context Learning Work? \textit{EMNLP 2022}.

\bibitem[Molnar, 2022]{molnar2022}
Molnar, C. (2022). \textit{Interpretable Machine Learning} (2nd ed.). Online Book.

\bibitem[Ouyang et al., 2022]{ouyang2022}
Ouyang, L., et al. (2022). Training Language Models to Follow Instructions with Human Feedback. \textit{NeurIPS 2022}.

\bibitem[Reiter \& Dale, 2000]{reiter2000}
Reiter, E., \& Dale, R. (2000). \textit{Building Natural Language Generation Systems}. Cambridge University Press.

\bibitem[Saltelli et al., 2008]{saltelli2008}
Saltelli, A., et al. (2008). \textit{Global Sensitivity Analysis: The Primer}. Wiley.

\bibitem[Schick et al., 2023]{schick2023}
Schick, T., et al. (2023). Toolformer: Language Models Can Teach Themselves to Use Tools. \textit{NeurIPS 2023}.

\bibitem[Vaswani et al., 2017]{vaswani2017}
Vaswani, A., et al. (2017). Attention Is All You Need. \textit{NeurIPS 2017}.

\end{thebibliography}

\end{document}
